
% Table created by stargazer v.5.2.3 by Marek Hlavac, Social Policy Institute. E-mail: marek.hlavac at gmail.com
% Date and time: Wed, May 29, 2024 - 10:33:01
\begin{table}[!htbp] \centering 
  \caption{Bank Credit Predicting Capital Excess Returns} 
  \label{bank credit and excess returns} 
\begin{tabular}{@{\extracolsep{10pt}}lcccc} 
\\[-1.8ex]\hline 
\hline \\[-1.8ex] 
 & \multicolumn{4}{c}{\textit{Dependent variable:}} \\ 
\cline{2-5} 
\\[-1.8ex] & \multicolumn{4}{c}{Average realized excess return $_{t+1}$} \\ 
\\[-1.8ex] & (1) Static Belief & (2) Bayesian & (3) Diagnostic  & (4) Data \\
\hline \\[-1.8ex] 
 $(\frac{\text{bank credit}}{\text{GDP}})_t$ & $-$0.01 & $-$0.01 & $-$0.02 & $-$0.02 \\ 
  &  &  &  &  \\ 
 \hline \\[-1.8ex] 
Observations &  &  &  & 867 \\ 
\hline 
\hline \\[-1.8ex] 
\multicolumn{5}{p{0.9\textwidth}}{\footnotesize \textit{Note}: 
  Model excess return is defined as the return to capital minus the risk-free rate. Data excess return is the excess 
equity index return from Online Appendix Table 3 of Baron and Xiong (2017). To ensure comparability, the model return 
to capital has been normalized to equal the standard deviation of returns reported by Baron and Xiong (2017).
}  
\end{tabular} 
\end{table} 
